\documentclass{article}
\usepackage[utf8]{inputenc}

\usepackage{a4}
\usepackage[margin=1in]{geometry}
\usepackage[doublespacing]{setspace}

\begin{document}

My name is Nir Elber, and I am a member of the number theory group of the mathematics department at MIT. In terms of timeline, I have just completed my first year, in which I took many classes; I have also chosen an advisor (Wei Zhang), and I will take my qualifying exam at the start of the fall semester. Roughly speaking, this means that I am in the final phases before starting research.

It may seem bizarre that a graduate student, who should essentially be a research student, has spent so much time not doing research. So let me say something about my studies. Number theory is traditionally interested in properties of the (positive) integers, such as polynomial equations they might satisfy. For some equations, it is easy to determine positive integer solutions, such as $x+y=z$; for others, a little more work is required, such as $x^2+y^2=z^2$. And for many, this task turns out to be exceedingly difficult, the poster child of which is positive integer solutions of $x^n+y^n=z^n$ for $n\ge3$. In general, it is computationally impossible to determine if an equation is ``easy'' or ``hard,'' but number theorists would like to try anyway. And so I have learned to sit down and apply the full force of modern mathematical machinery to these sorts of questions.

My focus in number theory lies in its intersection with algebraic geometry. The key observation here is that the ``shape'' of a polynomial equation informs number theory. (For example, one can take ``shape'' to mean how the set of solutions looks when it is graphed: the graphs of $y^2=x^2-x$ and $y^2=x^3-x$ look quite different!) Thus, one hopes to gain understanding in number theory by first having understanding in geometry. From my perspective, this means that I should study algebraic geometry to keep up with modern ideas.

Of course, there are many other fields of mathematics which have been turned into tools in service of number theory, and they roughly correspond to my coursework in my first year. Geometry is notoriously difficult, in part due to the fact that the underlying equations are non-linear. There are a variety of techniques to turn such ``non-linearities'' into ``linearities.'' In my opinion, representation theory is the most productive method, so I have also spent much effort devoted to it.
=
MIT is a school renowned for its faculty studying representation theory, which was something of a shock for me when I came from my undergraduate institution, UC Berkeley. If I were being reductive, I might describe UC Berkeley as a ``geometry school'' because so much of the number theory I consumed smelled like geometry. I owe a debt to Professors Martin Olsson and Yunqing Tang in particular. On the other hand, every other person I met at MIT studies representation theory, so peer pressure and a little imposter syndrome have pushed me further into this wonderful subject. I owe similar intellectual debts to Professors Roman Bezurkavnikov, Pavel Etingof, and Zhiwei Yun. (As is traditional in mathematics, names are listed alphabetically.)

I mentioned peer pressure and imposter syndrome as primary motivators. I am under the impression that these are a part of the universal graduate student experience, and they can easily destroy a career. However, the magic of MIT has made allowed me to appreciate these aspects as just another fuel once used carefully. Maybe I return from a reading group meeting with a few friends a little embarrassed because I made a few mistakes. That's fine---now I just have to spend a few hours correcting them on my own.

Such a mental adjustment takes effort and time, as have the many other adjustments I have experienced as I entered MIT. There are simply so many things to do, from seminars to classes, from academic meetings to social ones. I cannot imagine what my first year would have looked liked if I additionally had to teach, on top of all these other external and internal obligations, but I am fairly confident that I would have quickly spun out of control. So thank you for funding my first year. It has given me a strong foundation upon which to build the rest of my graduate and even academic carreer.

\end{document}